\chapter{Lecture 10}

\hfill {\footnotesize \textit{Date: 18/03/2026}}

\section{Tensor Product of Modules---continued}
\begin{proof}[Proof of existence]
    Consider \( M \times N \) simply as set. We form the free  \( R \)-module
     \[
        C \coloneq R^{(M \times N)},
    \]
     on this set, that is, \( C \) is defined to be the  \( R \)-module with a basis of formal symbols  \( (m, n) \in M \times N \) (we declare these new symbols to be linearly independent by definition).  So if  \( y \in C \), then \( y \) is a formal finite sum of the form  \( \sum a_i(m_i, n_i) \), where  \( a_i \in R, m_i \in M, \) and  \( n_i \in N \).

    Let \( D \) be a submodule of  \( C \) generated by all the elements of \( C \) of the following types:
    \[
        (m + m', n) - (m, n) - (m', n), (m, n + n') - (m, n) - (m, n'),
    \]
    \[
        (am, n) - a(m, n) \text{ and } (m, an) - a(m, n).
    \]
    Define the quotient module \( T \coloneq C / D \). We write the coset \( (m, n) + D \) in \( C / D \) as  \( m \otimes n \). So, elements of \( D \) are equivalent to \( 0 \) in  \( T = C / D \). So we get relations in  \( T = C / D \) like
    \[
        (m + m', n) \cong (m, n) + (m', n) \pmod{D},
    \]
    which is same as writing
    \[
        (m + m') \otimes n = m \otimes n + m' \otimes n.
    \]
    Similarly, we have
    \begin{gather*}
        m \otimes (n + n') = m \otimes n + m \otimes n',\\
        am \otimes n = a(m \otimes n) = m \otimes an.
    \end{gather*}
    These relations then gives a \( R \)-bilinear map \( g: M \times N \to T \), given by \( (m, n) \mapsto m \otimes n \).

    Given an \( R \)-bilinear map  \( f: M \times N \to P \), we want to show there exists a unique  \( R \)-linear map  \( \widetilde{f}: T \to P \) such that \( \widetilde{f} \circ g = f \).

    First define \( \overline{f}: C \to P \) by
    \[
        \overline{f}\left(\sum a_i(m_i, n_i)\right) = \sum a_i f(m_i, n_i).
    \]
    \( \overline{f} \) vanishes on all the elements of the submodule \( D \). Since \( f \) is an  \( R \)-bilinear map,
     \begin{gather*}
        \overline{f}((m + m', n) - (m, n) - (m', n)) = \overline{f}(m + m', n) - \overline{f}(m, n) - \overline{f}(m', n) = 0,\\
        \overline{f}((m, n + n') - (m, n) - (m, n')) = \overline{f}(m, n + n') - \overline{f}(m, n) - \overline{f}(m, n') = 0,\\
        \overline{f}((am, n) - a(m,n)) = \overline{f}(am, n) - \overline{f}(a(m,n)) = 0, \text{ and}\\
    \overline{f}((m,an) - a(m, n)) = \overline{f}(am, n) - \overline{f}a(m, n) = 0.
    \end{gather*}
    Since \( \overline{f} \) vanishes on the generators of \( D \), hence,  \( \overline{f} \) vanishes on all of \( D \). Therefore,  \( D \subseteq \ker \overline{f}\). So we can define the function \( \widetilde{f}: C / D \to P \) by
     \[
        \widetilde{f}(m \otimes n) = \overline{f}(m,n),
    \]
    and extend linearly to all of \( C / D \).  Then \( \widetilde{f} \) is  \( R \)-linear by construction. Now, we need to check if  \( \widetilde{f} \circ g = f \). It is enough to show this  \( \widetilde{f} \circ g = f \) for all  \( m\otimes n \in T \). And that is indeed true:
    \[
        \widetilde{f}(g(m, n)) =  \widetilde{f} (m \otimes n) = \overline{f}(m, n) = f(m,n).
    \]
    Now, if \( \widetilde{f}' \) is another map satisfying this property, that is,  \( \widetilde{f}' \circ g = f \), then for all generator \( m \otimes n \in T\),
     \begin{align*}
         & (\widetilde{f}' \circ g)(m, n) = (\widetilde{f} \circ g)(m, n)\\
         \implies &\widetilde{f}'(m \otimes n) = \widetilde{f}(m \otimes n)\\
    ,\end{align*}
    holds. So, \( \widetilde{f} = \widetilde{f}' \). So  \( \widetilde{f} \) is unique.
\end{proof}
\begin{example}
    Let \( M = \mathbb{Z}/m\mathbb{Z} \)  and \( N = \mathbb{Z}/n\mathbb{Z} \). We want to compute:
     \[
        \mathbb{Z}/m\mathbb{Z} \otimes_\mathbb{Z} \mathbb{Z}/n\mathbb{Z}.
     \]
    Since, \( \mathbb{Z}/m\mathbb{Z} \)
\end{example}
\todo{Notes for this lecture is incomplete.}
